% Options for packages loaded elsewhere
\PassOptionsToPackage{unicode}{hyperref}
\PassOptionsToPackage{hyphens}{url}
\PassOptionsToPackage{dvipsnames,svgnames,x11names}{xcolor}
\documentclass[
  10pt,
  fontset=fandol]{ctexart}
\usepackage{xcolor}
\usepackage[margin=16mm]{geometry}
\usepackage{amsmath,amssymb}
\setcounter{secnumdepth}{-\maxdimen} % remove section numbering
\usepackage{iftex}
\ifPDFTeX
  \usepackage[T1]{fontenc}
  \usepackage[utf8]{inputenc}
  \usepackage{textcomp} % provide euro and other symbols
\else % if luatex or xetex
  \usepackage{unicode-math} % this also loads fontspec
  \defaultfontfeatures{Scale=MatchLowercase}
  \defaultfontfeatures[\rmfamily]{Ligatures=TeX,Scale=1}
\fi
\usepackage{lmodern}
\ifPDFTeX\else
  % xetex/luatex font selection
\fi
% Use upquote if available, for straight quotes in verbatim environments
\IfFileExists{upquote.sty}{\usepackage{upquote}}{}
\IfFileExists{microtype.sty}{% use microtype if available
  \usepackage[]{microtype}
  \UseMicrotypeSet[protrusion]{basicmath} % disable protrusion for tt fonts
}{}
\makeatletter
\@ifundefined{KOMAClassName}{% if non-KOMA class
  \IfFileExists{parskip.sty}{%
    \usepackage{parskip}
  }{% else
    \setlength{\parindent}{0pt}
    \setlength{\parskip}{6pt plus 2pt minus 1pt}}
}{% if KOMA class
  \KOMAoptions{parskip=half}}
\makeatother
\setlength{\emergencystretch}{3em} % prevent overfull lines
\providecommand{\tightlist}{%
  \setlength{\itemsep}{0pt}\setlength{\parskip}{0pt}}
\usepackage{amsmath,amssymb,mathtools,needspace}
\xeCJKDeclareCharClass{CJK}{"2032,"0400->"04FF,"2160->"216B,"2460->"2473,"25A0,"25C7,"25CB,"25CF}
\IfFontExistsTF{Arial Unicode MS}{\xeCJKsetup{AutoFallBack=true}\setCJKfallbackfamilyfont{\CJKrmdefault}{Arial Unicode MS}}{}
\setcounter{secnumdepth}{0}
\setlength{\emergencystretch}{2em}
\allowdisplaybreaks
\usepackage{bookmark}
\IfFileExists{xurl.sty}{\usepackage{xurl}}{} % add URL line breaks if available
\urlstyle{same}
\hypersetup{
  pdftitle={带权的 Gauss 公式},
  colorlinks=true,
  linkcolor={Maroon},
  filecolor={Maroon},
  citecolor={Blue},
  urlcolor={Blue},
  pdfcreator={LaTeX via pandoc}}

\title{带权的 Gauss 公式}
\author{}
\date{}

\begin{document}
\maketitle

\subsubsection{4.4.5 带权的 Gauss 公式}\label{ux5e26ux6743ux7684-gauss-ux516cux5f0f}

考察积分 \(\displaystyle I=\int_a^b\rho(x)f(x)\,\mathrm{d}x\)，这里 \(\rho(x)\geq0\) 称为权函数，当 \(\rho(x)\equiv1\) 时即为普通的积分。

可以仿照处理普通积分的方法讨论带权的积分。考察求积公式

\[
\int_a^b\rho(x)f(x)\,\mathrm{d}x\approx\sum_{k=0}^{n}A_kf(x_k),
\]

如果它对于任意次数不超过 \(2n+1\) 的多项式均能准确地成立，则称之为 Gauss 型的。上述 Gauss 公式的求积节点 \(x_k\) 仍称为 Gauss 点。同样地，\(x_k\)（\(k=0,1,\cdots,n\)）是 Gauss 点的充要条件为，\(\displaystyle\omega(x)=\prod_{k=0}^{n}(x-x_k)\) 是区间 \([a,b]\) 上关于权函数 \(\rho(x)\) 的正交多项式。

若 \(a=-1,b=1\)，且取权函数 \(\rho(x)=\dfrac1{\sqrt{1-x^2}}\)，则所建立的 Gauss 公式为

\[
\int_{-1}^{1}\frac{f(x)}{\sqrt{1-x^2}}\,\mathrm{d}x
\approx\sum_{k=0}^{n}A_kf(x_k).
\tag{4.4.6}
\]

式（4.4.6）称为 Gauss-Chebyshev 公式。由于区间 \([-1,1]\) 上关于权函数 \(\dfrac1{\sqrt{1-x^2}}\) 的正交多项式是 Chebyshev 多项式（见 3.4.2 节），因此求积公式（4.4.6）的 Gauss 点是 \(n+1\) 次 Chebyshev 多项式的零点，即为

\[
x_k=\cos\left(\frac{2k+1}{2n+2}\pi\right)\qquad(k=0,1,\cdots,n).
\]

值得指出的是，运用正交多项式的零点构造 Gauss 求积公式，这种方法只是针对某些特殊的权函数才有效。构造 Gauss 公式的一般方法则是 4.1.2 节介绍过的待定系数法。

例如，设要构造下列形式的 Gauss 公式：

（公式续下页。）

\href{../../source-images/pdf-111.jpeg}{原页图像}

\subsubsection{4.4.5 带权的 Gauss 公式（续）}\label{ux5e26ux6743ux7684-gauss-ux516cux5f0fux7eed}

\[
\int_0^1\sqrt{x}f(x)\,\mathrm{d}x\approx A_0f(x_0)+A_1f(x_1),
\tag{4.4.7}
\]

令它对于 \(f(x)=1,x,x^2,x^3\) 准确成立，得

\[
\begin{cases}
A_0+A_1=\dfrac23,\\[5pt]
x_0A_0+x_1A_1=\dfrac25,\\[5pt]
x_0^2A_0+x_1^2A_1=\dfrac27,\\[5pt]
x_0^3A_0+x_1^3A_1=\dfrac29.
\end{cases}
\tag{4.4.8}
\]

由于

\[
x_0A_0+x_1A_1=x_0(A_0+A_1)+(x_1-x_0)A_1,
\]

利用式（4.4.8）的第一式，可将第二式化为

\[
\frac23x_0+(x_1-x_0)A_1=\frac25.
\]

同样地，利用第二式化第三式，利用第三式化第四式，分别得

\[
\frac25x_0+(x_1-x_0)x_1A_1=\frac27,\qquad
\frac27x_0+(x_1-x_0)x_1^2A_1=\frac29.
\]

从上面三个式子中消去 \((x_1-x_0)A_1\)，有

\[
\begin{cases}
\dfrac25x_0+\left(\dfrac25-\dfrac23x_0\right)x_1=\dfrac27,\\[5pt]
\dfrac27x_0+\left(\dfrac27-\dfrac25x_0\right)x_1^2=\dfrac29,
\end{cases}
\]

即

\[
\begin{cases}
\dfrac25(x_0+x_1)-\dfrac23x_0x_1=\dfrac27,\\[5pt]
\dfrac27(x_0+x_1)-\dfrac25x_0x_1=\dfrac29.
\end{cases}
\]

由此解出

\[
x_0x_1=\frac5{21},\quad x_0+x_1=\frac{10}9,
\]

从而求出

\[
\begin{aligned}
x_0&=0.821\,162,\quad x_1=0.289\,949,\\
A_0&=0.389\,111,\quad A_1=0.277\,556.
\end{aligned}
\]

于是形如式（4.4.7）的 Gauss 公式是

\[
\int_0^1\sqrt{x}f(x)\,\mathrm{d}x
\approx0.389\,111f(0.821\,162)+0.277\,556f(0.289\,949).
\]

\begin{quote}
校注（agent 补充，CH04B-E003）：原扫描在``从上面三个式子中消去''后的第二行印为 \(\dfrac27x_0+\left(\dfrac27-\dfrac25x_0\right)x_1^2=\dfrac29\)，本转录原样保留。疑误是括号后 \(x_1^2\) 的指数：前一行有 \((x_1-x_0)x_1A_1=\dfrac27-\dfrac25x_0\)，代入再上一组的第四式应只再乘 \(x_1\)，故此处应为 \(\left(\dfrac27-\dfrac25x_0\right)x_1\)，才与紧接着``即''后的第二行相符。这是原书疑似排印错误，非 OCR 错误。\href{../assets/ch04b/review-p111-elimination.png}{局部原扫描}。
\end{quote}

\subsection{本节覆盖原页的校注记录}\label{ux672cux8282ux8986ux76d6ux539fux9875ux7684ux6821ux6ce8ux8bb0ux5f55}

下列记录按原页关联，含同页相邻小节的校注；计算时同时读取上文校注和完整原页。

\begin{itemize}
\tightlist
\item
  \texttt{CH04B-E003}：原书 PDF 页 111
\end{itemize}

\end{document}
